\documentclass[11pt,a4paper]{article}

% ============================================
% ENCODING & FONTS
% ============================================
\usepackage[utf8]{inputenc}
\usepackage[T1]{fontenc}
\usepackage{lmodern}

% ============================================
% PAGE LAYOUT
% ============================================
\usepackage[margin=1.618cm, top=2.618cm, bottom=2.618cm]{geometry}

% ============================================
% ESSENTIAL PACKAGES
% ============================================
\usepackage{float}
\usepackage{caption}
\usepackage{setspace}
\usepackage{fancyhdr}
\usepackage{xcolor}
\usepackage{hyperref}
\usepackage{amsmath}
\usepackage{amssymb}
\usepackage{amsthm}
\usepackage{booktabs}
\usepackage{array}
\usepackage{graphicx}
\usepackage{longtable}
\ifdefined\DeclareUnicodeCharacter
  \DeclareUnicodeCharacter{00B0}{\ensuremath{^\circ}}
\fi

% ============================================
% TITLE FORMATTING
% ============================================
\usepackage{titling}
\pretitle{\LARGE\bfseries}
\posttitle{\vspace{-0.4em}}
\preauthor{}
\postauthor{}
\predate{}
\postdate{}
\setlength{\droptitle}{-2.0em}

% ============================================
% HEADER/FOOTER
% ============================================
\setlength{\headheight}{14pt}
\pagestyle{fancy}
\fancyhf{}
\fancyhead[L]{A Certified Torsion-Free $\mathrm{G}_2$ Structure on a TCS Neck Model}
\fancyhead[R]{\thepage}
\renewcommand{\headrulewidth}{0.2pt}

% ============================================
% HYPERREF
% ============================================
\hypersetup{
  colorlinks=true,
  linkcolor=blue,
  citecolor=blue,
  urlcolor=blue,
  pdftitle={A Certified Torsion-Free G2 Structure on a TCS Neck Model via Computer-Assisted Proof},
  pdfauthor={Brieuc de La Fourniere}
}

% ============================================
% SPACING
% ============================================
\setstretch{1.2}
\setlength{\parskip}{0.4em}
\setlength{\parindent}{0pt}

% ============================================
% THEOREM ENVIRONMENTS
% ============================================
\newtheorem{theorem}{Theorem}[section]
\newtheorem{proposition}[theorem]{Proposition}
\newtheorem{lemma}[theorem]{Lemma}
\newtheorem{corollary}[theorem]{Corollary}
\theoremstyle{definition}
\newtheorem{definition}[theorem]{Definition}
\theoremstyle{remark}
\newtheorem{remark}[theorem]{Remark}

% ============================================
% CUSTOM COMMANDS
% ============================================
\newcommand{\eps}{\varepsilon}
\newcommand{\CIvvv}{\mathrm{CI}(2,2,2)}
\newcommand{\CIvvvv}{\mathrm{CI}(1,2,2,2)}
\newcommand{\Psix}{\mathbb{P}^6}
\newcommand{\RR}{\mathbb{R}}
\newcommand{\Gtwo}{\mathrm{G}_2}
\newcommand{\Vol}{\mathrm{Vol}}
\newcommand{\bvec}[1]{\boldsymbol{#1}}
\newcommand{\norm}[1]{\left\|#1\right\|}

\pdfstringdefDisableCommands{%
  \def\CIvvv{CI(2,2,2)}%
  \def\CIvvvv{CI(1,2,2,2)}%
  \def\Psix{P^6}%
  \def\Gtwo{G2}%
}

% ============================================
% DOCUMENT
% ============================================

\title{%
\LARGE\textbf{A Certified Torsion-Free $\Gtwo$ Structure\\[0.2em]
on a TCS Neck Model via Computer-Assisted Proof}
}
\author{}
\date{}

\begin{document}

\maketitle
\noindent\rule{\textwidth}{0.2pt}

\noindent\textbf{Author}: Brieuc de La Fourni\`ere

\noindent Independent researcher, Beaune, France

\noindent Contact: \href{mailto:brieuc@bddelaf.com}{\texttt{brieuc@bddelaf.com}}

\noindent\textbf{Date}: April 2026

\noindent\textbf{DOI}: \href{https://doi.org/10.5281/zenodo.19892350}{10.5281/zenodo.19892350}

\vfill

\begin{abstract}
We construct an explicit $\Gtwo$ structure on the TCS-type neck model $K3 \times T^2 \times [0,1]$ and certify, via Newton--Kantorovich interval arithmetic with zero finite differences, the existence of a unique nearby torsion-free $\Gtwo$ structure with contraction parameter $h \leq 8.95 \times 10^{-9}$ (analytical certificate, $\beta = 0.321$ from a Sturm--Liouville bound), sharpened to $h \leq 1.43 \times 10^{-9}$ with the numerical $\beta = 0.02961$; both well below the threshold $1/2$ and relative to the fixed K3 neural-network input. The solution space is finite-dimensional ($\RR^{168}$), the torsion map is a polynomial of degree $\leq 10$ in the neck coordinate, and the analytical inverse bound is derived from the algebraically exact determinant $\det(g) = 65/32$.

The optimized 169 parameters (168 Chebyshev--Cholesky coefficients $+$ ACyl decay rate $\gamma$) exhibit a five-constant approximate structure ($g_{ss} \approx 19/6$, $g_{T^2} \approx 7/6$, $g_{K3} \approx 64/77$, $\det(g) = 65/32$, $\gamma = 2\pi\sqrt{6/7}$) with sub-percent deviations; the first two are torsion minimizers and the last is derived from the $T^2$ Hodge Laplacian and $H^1(K3) = 0$. The $\Gtwo$ torsion-free condition acts as an eigenvalue equalizer, reducing the K3 fiber anisotropy from 10:1 to 1.012:1. We prove analytically that product-type Ricci-flat $\Gtwo$ metrics satisfy spectral democracy (Theorem~\ref{thm:D}).

The NK certificate extends to the full 7D product metric via an analytical Fr\'echet bound on the K3 fiber, giving $\|T(g^*)\|_{C^0} \leq 1.59 \times 10^{-3}$, well below the Joyce perturbation threshold $\varepsilon_0 \approx 0.1$. The underlying seam-sector torsion at the NK iterate is $\eta \leq 2.949 \times 10^{-5}$ (direct sup of $\sqrt{|d\varphi|^2 + |d\star\varphi|^2}$ at Clenshaw--Curtis nodes, the value used in NK; the slightly larger $2.984 \times 10^{-5}$ in Table~\ref{tab:torsion_iter} is the conservative $\sqrt{(\sup|d\varphi|)^2 + (\sup|d\star\varphi|)^2}$), and the empirical 7D maximum measured directly over $220{,}000$ K3 fiber points is $3.154 \times 10^{-5}$. Conditional on the existence of a compact simply connected 7-manifold with Betti numbers $(b_2, b_3) = (21, 77)$ (a pair absent from all known $\Gtwo$ constructions), this certificate yields a torsion-free $\Gtwo$ metric on the compact manifold via Joyce's perturbation theorem.
\end{abstract}

\noindent\textbf{Keywords}: $\Gtwo$ holonomy; computer-assisted proof; Newton--Kantorovich; interval arithmetic; Chebyshev--Cholesky metric; spectral democracy.

\vfill
\noindent\rule{\textwidth}{0.2pt}

\newpage
\tableofcontents

\newpage

% ============================================
\section{Introduction}
% ============================================

\subsection{Context}

A compact Riemannian 7-manifold $(M, g)$ with holonomy $\mathrm{Hol}(g) = \Gtwo$ carries a torsion-free $\Gtwo$-structure: a 3-form $\varphi$ satisfying $d\varphi = 0$ and $d\star\varphi = 0$ \cite{Joyce2000,Bryant1987}. The metric is determined by $\varphi$, and torsion-freeness implies Ricci-flatness. Such manifolds play a central role in geometric analysis and in M-theory compactifications \cite{AcharyaWitten2001,Acharya1999}.

Despite significant progress on existence results (through Joyce's orbifold resolutions \cite{Joyce2000} and the Twisted Connected Sum (TCS) construction of Kovalev \cite{Kovalev2003} and Corti--Haskins--Nordstr\"om--Pacini \cite{CHNP2015}), explicit metric data on compact $\Gtwo$ manifolds has remained unavailable. The known constructions are perturbative: they establish existence of torsion-free metrics near approximate ones, without producing numerical metric tensor components. This contrasts with the Calabi--Yau setting, where numerical metrics have been computed via Donaldson's algorithm \cite{Donaldson2009}, balanced metrics \cite{HeadrickNassar2013}, and neural network methods \cite{LarforsLukasRuehle2022,Anderson2021}.

\subsection{Main results}

In this paper, we address part of this gap by constructing an explicit, NK-certified $\Gtwo$ structure on the TCS-type neck model $K3 \times T^2 \times [0,1]$, with seam-sector geometry adapted to a putative compact 7-manifold of Betti numbers $(b_2, b_3) = (21, 77)$ (§§\ref{sec:metric}--\ref{sec:torsion}); we certify its torsion-free completion via Newton--Kantorovich contraction (§\ref{sec:nk}) and prove an analytical spectral democracy theorem (§\ref{sec:democracy}). The certificate concerns the neck model itself; its extension to a compact realization, via Joyce's perturbation theorem, is conditional on the existence of such a manifold (§\ref{sec:manifold}). Spectral data (gap, Betti confirmation, harmonic forms) are presented in a companion paper [B].

Our main results are:

\begin{proposition}[Metric certification]\label{prop:A}
There exists a $\Gtwo$-structure $\varphi_0$, constructed as a product-type metric $g = g_{\mathrm{seam}}(s) \otimes g_{T^2} \otimes g_{K3}(y)$ on a TCS-type neck, with the following certified properties:

\emph{(i)} The full 7D torsion satisfies $\|T(\varphi_0)\|_{C^0} \leq 1.59 \times 10^{-3}$ (analytical Fr\'echet-certified upper bound, §\ref{sec:frech}); the empirical value measured directly over $220{,}000$ K3 fiber points is much tighter, $\|T(\varphi_0)\|_{C^0,\mathrm{emp}} \leq 3.154 \times 10^{-5}$.

\emph{(ii)} The Newton--Kantorovich contraction parameter $h = \beta\eta\omega = 1.43 \times 10^{-9} < 1/2$ (margin $\times 350$ million), certifying existence of a unique torsion-free $\Gtwo$-structure $\varphi^*$ with $\mathrm{dist}(\varphi^*, \varphi_0) \leq 8.73 \times 10^{-7}$, $\delta g/g \leq 1.35 \times 10^{-7}$.

\emph{(iii)} The metric is positive definite with $\det(g) = 65/32$ (algebraically exact), $\kappa(g) \leq 3.88$, and $\lambda_{\min}(g) \geq 0.817$.
\end{proposition}

\begin{corollary}[Metric decomposition]
The 169 parameters (168 Chebyshev--Cholesky $+$ $\gamma$) reduce to 5 structural constants ($g_{ss} = 19/6$, $g_{T^2} = 7/6$, $g_{K3} = 64/77$ (mean), $\det(g) = 65/32$, $\gamma = 2\pi\sqrt{6/7}$) with sub-percent corrections. Even Chebyshev modes vanish to machine precision; odd modes are Cholesky gauge.
\end{corollary}

\begin{theorem}[Spectral democracy]\label{thm:D}
Let $(N, g)$ be a Riemannian manifold with product-type metric $g = g_B(s) \oplus g_F$ where $g_F$ is independent of $s$ and $\mathrm{Ric}(g) = 0$. For a transverse 1-form $\alpha = f(s)\,d\theta$ with $\theta$ a fiber coordinate and $g_{\theta\theta}$ constant:
\[
\Delta_1(f\, d\theta) = (\Delta_0 f)\, d\theta.
\]
In particular, the transverse 1-form and scalar Laplacian spectra coincide.
\end{theorem}

\subsection{Status of results}

\begin{table}[ht]
\centering
\begin{tabular}{ll}
\toprule
Level & Definition \\
\midrule
\textbf{Algebraic} & Holds rigorously by construction (e.g., Cholesky SPD, determinant normalization) \\
\textbf{Analytically proved} & Derived from first principles with complete proof \\
\textbf{Certified} & Established by computer-assisted proof with interval arithmetic \\
\textbf{Numerical} & Computed with controlled discretization error and convergence analysis \\
\textbf{Observed} & Empirical match validated by independent methods but not derived \\
\textbf{Conditional} & Depends on the topological realization of a compact $M$ with $(b_2, b_3) = (21,77)$ \\
\bottomrule
\end{tabular}
\caption{Epistemological levels used in this paper.}
\end{table}

\begin{table}[ht]
\centering
\small
\begin{tabular}{lll}
\toprule
Result & Level & Depends on \\
\midrule
Prop.~\ref{prop:A}: NK existence of torsion-free metric & Certified & Banach space framework (§\ref{sec:nk}) \\
Prop.~\ref{prop:A}: $\mathrm{Hol}(g^*) = \Gtwo$ & Follows from Joyce \cite{Joyce2000} & Prop.~\ref{prop:A} + simple connectivity \\
\textbf{Thm.~\ref{thm:D}: Spectral democracy} & \textbf{Analytically proved} & Ricci-flatness + product metric \\
Corollary: $169 \to 5$ constants & Observed, Lean-verified & Prop.~\ref{prop:A} \\
Prop.~\ref{prop:E}: Adiabatic decomposition & Certified & NK proximity + product structure \\
Structural identities (§\ref{sec:decomp}) & Observed, validated & Torsion minimization + Lean \\
All results on $M$ & Conditional & Compact $M$ with $(21, 77)$; open (§\ref{sec:manifold}) \\
\bottomrule
\end{tabular}
\caption{Summary of results and their epistemological status.}
\end{table}

\subsection{Relation to existing work}

The closest precedent for numerical metrics on special-holonomy manifolds is Donaldson's algorithm for Calabi--Yau metrics \cite{Donaldson2009}, subsequently refined by Headrick--Nassar \cite{HeadrickNassar2013} for K3 surfaces and extended by Larfors--Lukas--Ruehle \cite{LarforsLukasRuehle2022} and Anderson et al.\ \cite{Anderson2021} using neural networks. In the $\Gtwo$ setting, Heyes--Hirst--S\'a Earp--Silva \cite{HeyesHirst2026} recently computed neural $\Gtwo$-structures on non-compact cohomogeneity-one examples. Our work extends these to the compact case with NK certification.

Previous numerical work on $\Gtwo$ metrics has focused on non-compact examples. For compact $\Gtwo$ manifolds, the existence results of Joyce \cite{Joyce2000,JoyceII1996} and the TCS construction \cite{Kovalev2003,CHNP2015} establish metrics perturbatively but do not yield explicit numerical values. The present paper combines the construction, certification, and analytical decomposition in a single self-contained treatment.

An independent arithmetic derivation of $(b_2, b_3) = (21, 77)$ has recently been proposed by Zhou \& Zhou~\cite{ZhouZhouA2026,ZhouZhouE2026}: combining spectral self-referential dynamics with Diophantine constraints arising from anomaly-cancellation screening principles on $\Gtwo$ manifolds with ADE singularities, they show that $(21, 77)$ is the unique positive integer solution of $b_2 + b_3 = 98$, $11 b_2 = 3 b_3$, subject to three arithmetic screening criteria. Their approach is complementary to the present one: the present paper constructs an explicit metric realizing $(21, 77)$ and certifies its $\Gtwo$ structure; \cite{ZhouZhouA2026,ZhouZhouE2026} argue from algebraic stability principles that $(21, 77)$ is uniquely distinguished.

\subsection{Outline}

Section~\ref{sec:manifold} discusses the topological context and the status of $(21, 77)$ in the $\Gtwo$ landscape. Section~\ref{sec:metric} presents the full 7D metric construction: the three adiabatic sectors, coordinate conventions, the Chebyshev--Cholesky parametrization, the reconstruction algorithm, the complete metric decomposition with observed rational approximations (Proposition~\ref{prop:F}) and torsion-minimizer validation (Proposition~\ref{prop:G}), and the adiabatic decomposition lemma (Proposition~\ref{prop:E}) establishing certified spectral control on the torsion-free metric. Section~\ref{sec:torsion} analyzes the torsion, including the $\Gtwo$ representation-theoretic structure and K3 fiber verification. Section~\ref{sec:nk} presents the NK certification in detail. Section~\ref{sec:democracy} proves the spectral democracy theorem analytically. Section~\ref{sec:discussion} discusses limitations and open questions.

% ============================================
\section{The Manifold and its Topological Context}\label{sec:manifold}
% ============================================

\subsection{Setup}

We construct an NK-certified $\Gtwo$ metric on a 7-manifold model whose Hodge Laplacian yields $b_2 = 21$ harmonic 2-forms and $b_3 = 77$ harmonic 3-forms (companion paper [B]). If a compact, simply connected 7-manifold $M$ with Betti numbers $(b_2, b_3) = (21, 77)$ and holonomy $\Gtwo$ exists, its full Betti spectrum is determined by Poincar\'e duality: $(1, 0, 21, 77, 77, 21, 0, 1)$ with Euler characteristic $\chi(M) = 0$.

The NK certificate (Proposition~\ref{prop:A}) and the spectral democracy theorem (Theorem~\ref{thm:D}) are properties of the explicit certified metric, independent of any global compactification hypothesis.

\subsection{The pair $(21, 77)$ in the landscape of known $\Gtwo$ manifolds}

The pair $(b_2, b_3) = (21, 77)$ does not appear among previously constructed compact $\Gtwo$ manifolds.

\textbf{Twisted connected sums (TCS).} The Kovalev--CHNP construction \cite{Kovalev2003,CHNP2015} glues two asymptotically cylindrical Calabi--Yau threefolds along a neck region $K3 \times T^2 \times I$. The CHNP 2015 tabulation covers building blocks with Picard rank $\rho \leq 9$, yielding $b_2 \leq 18$. For orthogonal gluing, Lemma 6.7 of \cite{CHNP2015} implies $b_2 + b_3$ is always odd; since $21 + 77 = 98$ is even, orthogonal TCS is excluded. Non-orthogonal TCS or extra-twisted connected sums \cite{Nordstrom2023} do not have this parity constraint and remain open.

\textbf{Joyce orbifold resolutions.} Joyce \cite{Joyce2000} constructs compact $\Gtwo$ manifolds by resolving singularities of flat orbifolds $T^7/\Gamma$, obtaining 252 distinct Betti number pairs. The pair $(21, 77)$ lies within the bounding box but does not appear among the 252 documented types.

\textbf{Topological existence.} By Wilkens' classification \cite{Wilkens1972}, any pair $(b_3, p_1)$ with $p_1 \equiv 0 \pmod{4}$ is realized by a smooth closed 2-connected 7-manifold, guaranteeing the existence of a smooth 7-manifold with $b_3 = 77$ as a topological space, but not a $\Gtwo$-holonomy metric.

\textbf{Summary.} The pair $(21, 77)$ is compatible with all known constraints on compact $\Gtwo$ manifolds but has not yet been realized by any explicit construction. A complete geometric construction remains an open problem.

\subsection{Metric model and assumptions}

The metric is constructed on a local model with coordinates $(s, \theta, \chi, y_1, y_2, y_3, y_4)$ adapted to a neck region diffeomorphic to $K3 \times T^2 \times [0, 1]$. The metric extends to $s \in [-2, 3]$ via exponential ACyl decay (§\ref{sec:acyl}) with rate $\gamma = 2\pi\sqrt{6/7} \approx 5.817$.

For the holonomy argument (§\ref{sec:holonomy}), we note that a torsion-free $\Gtwo$-structure on a compact simply connected manifold has holonomy exactly $\Gtwo$ \cite[Prop.~11.2.3]{Joyce2000}.

% ============================================
\section{The Metric}\label{sec:metric}
% ============================================

\subsection{Setup and conventions}

Throughout this paper, $M$ denotes the 7-manifold model. We work in local coordinates $(s, \theta, \chi, y_1, y_2, y_3, y_4)$ adapted to a neck region diffeomorphic to $K3 \times T^2 \times [0, 1]$, where $s$ is the neck (seam) coordinate, $(\theta, \chi)$ parametrize $T^2$, and $(y_1, \ldots, y_4)$ are K3 fiber coordinates.

The $\Gtwo$-structure is specified by the associative 3-form $\varphi \in \Omega^3(M)$ and the coassociative 4-form $\psi = \star\varphi \in \Omega^4(M)$. Torsion-freeness requires $d\varphi = 0$ and $d\star\varphi = 0$.

\subsection{The full 7D metric}

The metric on the neck region decomposes into three adiabatic sectors:
\begin{equation}
g(s, \theta, \chi, y) = g_{\mathrm{seam}}(s) \oplus g_{T^2} \oplus g_{K3}(y).
\end{equation}

\textbf{(i) Seam sector} $g_{\mathrm{seam}}(s)$: a one-parameter family of $7 \times 7$ positive-definite matrices along $s$, parametrized by 169 Chebyshev--Cholesky coefficients (§\ref{sec:cheb}). This sector carries $99.93\%$ of the total torsion.

\textbf{(ii) Torus sector} $g_{T^2}$: flat metric on the $T^2$ fiber, isotropic to $|g^{\theta\theta} - g^{\chi\chi}| = 3 \times 10^{-7}$.

\textbf{(iii) K3 sector} $g_{K3}(y)$: Ricci-flat K\"ahler metric on the K3 fiber $\CIvvvv \subset \Psix$, computed via the cymyc neural-network library \cite{cymyc2024} (extending the framework of \cite{LarforsLukasRuehle2022}) on $220{,}000$ sample points with RMS accuracy $\sigma = 0.011$. \emph{Note}: $\CIvvvv \subset \Psix$ and $\CIvvv \subset \mathbb{P}^5$ denote the same K3 family; the linear factor in $\CIvvvv$ embeds the surface as a hyperplane section of $\Psix$, and restricting to that hyperplane recovers $\CIvvv \subset \mathbb{P}^5$. Period calculations and Picard--Fuchs analysis elsewhere in this paper use the $\CIvvv$ notation.

\textbf{Adiabatic dominance.} The product structure implies the seam sector dominates the torsion. This is a certified \emph{a posteriori} result (§\ref{sec:frech}):

\begin{table}[ht]
\centering
\small
\begin{tabular}{lll}
\toprule
Test & Result & Status \\
\midrule
K3 torsion contribution (Fr\'echet, 220k pts) & $0.07\%$ of total & Certified \\
$T^2$ isotropy & $3 \times 10^{-7}$ & Certified \\
Fiber Fourier $k=1$ extension & $0\%$ torsion reduction & Null experiment \\
KK gauge field extension & $0\%$ torsion reduction & Null experiment \\
K3 fiber metric variation along $s$ & $< 0.002\%$ & Certified \\
Mean torsion shift with K3 coupling & $0.000006\%$ & Certified ($2000 \times 100$ nodes) \\
\bottomrule
\end{tabular}
\caption{Adiabatic dominance verification tests.}
\end{table}

\subsection{Coordinate system}

We use an atlas with three charts: $U_L$ ($s \in [-2, 0]$), $U_N$ ($s \in [0, 1]$), $U_R$ ($s \in [1, 3]$). The fiber coordinate index convention at fixed $s$:

\begin{table}[ht]
\centering
\begin{tabular}{lll}
\toprule
Index & Direction & Type \\
\midrule
0 & $s$ (seam) & Neck parameter \\
1 & $\theta$ (circle fiber) & $T^2$ \\
2 & $y_1$ (K3 real part 1) & K3 \\
3 & $y_2$ (K3 imaginary part 1) & K3 \\
4 & $y_3$ (K3 real part 2) & K3 \\
5 & $y_4$ (K3 imaginary part 2) & K3 \\
6 & $\chi$ (circle fiber) & $T^2$ \\
\bottomrule
\end{tabular}
\caption{Fiber coordinate convention.}
\end{table}

\subsection{Chebyshev--Cholesky parametrization}\label{sec:cheb}

To ensure positive-definiteness, we write $g_{\mathrm{seam}}(s) = L(s)L(s)^T$ where $L$ is lower-triangular. The 28 independent entries of $L$ are expanded in Chebyshev polynomials of degree $K = 5$:
\begin{equation}
L_{\mathrm{flat},j}(s) = \sum_{k=0}^{K} c_{kj}\, T_k(2s - 1), \qquad j = 0, \ldots, 27.
\end{equation}

\textbf{Softplus activation on the diagonal.} The diagonal entries of $L$ are passed through a softplus function:
\begin{equation}
L_{ii} = \log(1 + \exp(L_{\mathrm{flat},j})) \quad \text{for } j \in \{0, 2, 5, 9, 14, 20, 27\}.
\end{equation}

\textbf{Determinant normalization.} At each $s$, the metric is rescaled to enforce $\det(g(s)) = 65/32$ (exactly, by construction) via the scaling $L(s) \to \alpha(s) \cdot L(s)$ where $\alpha^7 \det(L) = \sqrt{65/32}$.

\textbf{Total parameter count:} $6$ Chebyshev modes $\times$ 28 Cholesky entries $= 168$ parameters, plus 1 ACyl decay rate $\gamma = $ \textbf{169 parameters total}.

\subsection{ACyl extension}\label{sec:acyl}

For $s$ outside the neck $[0, 1]$, the metric decays exponentially toward the asymptotic Calabi--Yau cross-section metric:
\begin{equation}
g(s) \to g_\infty + [g(s_{\mathrm{bdy}}) - g_\infty] \cdot \exp(-2\gamma|s - s_{\mathrm{bdy}}|)
\end{equation}
with NK-computed decay rate $\gamma_{\mathrm{NK}} = 5.811$ (compare $\gamma = 2\pi\sqrt{6/7} \approx 5.817$ derived in §\ref{sec:decomp}; relative deviation $0.1\%$) and bulk domain extending to $s \in [-2, 3]$.

Since $H^1(K3) = 0$ (K3 is simply connected), the leading transverse 3-form modes lie in $\Omega^1(T^2) \otimes H^2(K3)$, giving $\gamma^2 = 4\pi^2/g_{T^2}$. With $g_{T^2} = 7/6$ this yields $\gamma = 2\pi\sqrt{6/7} \approx 5.817$; the NK-computed rate $\gamma_{\mathrm{NK}} = 5.811$ is $0.1\%$ below, consistent with the $0.2\%$ NK proximity of $g_{T^2}$.

\subsection{Metric properties}

At the neck midpoint $s = 0.5$, the metric is approximately block-diagonal:
\begin{table}[ht]
\centering
\begin{tabular}{lll}
\toprule
Direction & Block eigenvalue & Role \\
\midrule
Seam ($s$) & $3.166$ & Dominant \\
$T^2$ ($\theta, \chi$) & $1.169$ & Intermediate \\
K3 ($y_1\ldots y_4$) & $0.828$ & Fiber \\
\bottomrule
\end{tabular}
\caption{Block eigenvalues at $s = 0.5$ (cf.~\S\ref{sec:decomp}, Phase 1).}
\end{table}

Off-diagonal entries are below $0.007$. Certified bounds (Gershgorin, interval arithmetic): $\lambda_{\min}(g) \geq 0.817$, $\lambda_{\max}(g) \leq 3.167$, $\kappa(g) \leq 3.88$.

\begin{figure}[ht]
\centering
\includegraphics[width=\textwidth]{fig_metric_profile_paper_a.png}
\caption{Metric profile at the neck midpoint $s = 0.5$. \emph{(a)} The $7 \times 7$ tensor $g$ (det-normalised to $65/32$): nearly block-diagonal seam $|$ $T^2$ $|$ $4 \times K3$ structure with off-diagonal couplings below $0.007$. \emph{(b)} Block eigenvalues coloured by adiabatic sector, with rational approximations $g_{ss} \approx 19/6$, $g_{T^2} \approx 7/6$, $g_{K3} \approx 64/77$ marked, and the ACyl decay rate $\gamma = 2\pi\sqrt{6/7}$ shown for reference (NK-computed $\gamma_{\mathrm{NK}} = 5.811$, $0.1\%$ proximity).}
\label{fig:metric_profile}
\end{figure}

\subsection{Reconstruction algorithm}

Given the coefficient matrix $C = (c_{kj})$ and decay rate $\gamma$, the metric $g(s)$ at any point $s$ is reconstructed in 6 steps:
\begin{enumerate}
\item \textbf{Chebyshev evaluation:} For each $j \in \{0,\ldots,27\}$, compute $L_{\mathrm{flat},j}(s) = \sum_k c_{kj} T_k(2s-1)$.
\item \textbf{Reshape:} Pack the 28 values into a $7\times 7$ lower-triangular matrix $L(s)$.
\item \textbf{Softplus diagonal:} Replace $L_{ii} \leftarrow \log(1 + \exp(L_{ii}))$ for the 7 diagonal entries.
\item \textbf{Metric:} Compute $g(s) = L(s) \cdot L(s)^T$.
\item \textbf{Det-normalize:} Rescale $g(s) \leftarrow g(s) \cdot (65/32 / \det(g(s)))^{1/7}$.
\item \textbf{ACyl tails:} For $s \notin [0,1]$, apply exponential decay with rate $\gamma$.
\end{enumerate}

The $\Gtwo$ 3-form $\varphi$ and 4-form $\psi = \star\varphi$ are constructed from the Cholesky factor $L$ and the Fano plane structure constants:
\begin{equation}
\varphi_{ijk}(s) = \sum_{(abc) \in \mathrm{Fano}} \mathrm{sign}(abc) \cdot L_{ia}(s)\, L_{jb}(s)\, L_{kc}(s).
\end{equation}

\subsection{Complete metric decomposition}\label{sec:decomp}

The adiabatic product structure $g \approx g_{\mathrm{seam}} \oplus g_{T^2} \oplus g_{K3}$ admits a characterization in terms of topological invariants of the TCS construction.

\begin{proposition}[Observed rational approximations]\label{prop:F}
The three adiabatic sectors have effective eigenvalues approximated by irreducible rationals with deviations below $0.5\%$:
\begin{center}
\small
\begin{tabular}{lllll}
\toprule
Sector & Symbol & Rational match & NK numerical & Deviation \\
\midrule
Seam & $g_{ss}$ & $\frac{19}{6}$ & $3.16559$ & $0.03\%$ \\
$T^2$ fiber & $g_{T^2}$ & $\frac{7}{6}$ & $1.16900$ & $0.20\%$ \\
K3 fiber (mean) & $g_{K3}$ & $\frac{64}{77}$ & $0.8278$ & $0.41\%$ \\
\bottomrule
\end{tabular}
\end{center}
\end{proposition}

\begin{proposition}[Torsion minimum]\label{prop:G}
Forcing the exact topological fractions $g_{ss} = 19/6$ and $g_{T^2} = 7/6$ in the Chebyshev expansion \textbf{lowers} the mean torsion $\|\nabla\varphi\|$ by $0.05\%$. The topological values are not imposed; they emerge as torsion minimizers.
\end{proposition}

\textbf{ACyl decay.} The exponential matching rate $\gamma$ satisfies:
\begin{equation}
\gamma^2 = \frac{4\pi^2}{g_{T^2}} = \frac{24\pi^2}{7} \approx 33.839, \qquad \gamma = \frac{2\pi}{\sqrt{g_{T^2}}} = 2\pi\sqrt{\frac{6}{7}} \approx 5.817.
\end{equation}
The derivation uses $H^1(K3) = 0$: since K3 is simply connected, the first transverse 3-form eigenvalue on $T^2 \times K3$ is the first $T^2$ eigenvalue $\lambda_1(T^2) = 4\pi^2/g_{T^2}$ (K3 harmonic 2-forms contribute zero to the transverse spectrum). With $g_{T^2} = 7/6$ this gives $\gamma = 2\pi\sqrt{6/7}$. The NK-computed rate $\gamma_{\mathrm{NK}} = 5.811$ lies $0.1\%$ below, consistent with the $0.2\%$ NK proximity of $g_{T^2}$.

\textbf{Determinant consistency and the status of $g_{K3}$.} The metric determinant $\det(g) = 65/32$ is imposed exactly by construction. The rational triple $(g_{ss}, g_{T^2}, g_{K3}) = (19/6, 7/6, 64/77)$ is faithful to the measured block eigenvalues at $0.03\%$--$0.41\%$ precision, but the naive block-isotropic product $g_{ss} \cdot g_{T^2}^2 \cdot g_{K3}^4 \approx 2.0571$ differs from $\det(g) = 65/32 = 2.03125$ by $1.27\%$. The two rationals $g_{ss} = 19/6$ and $g_{T^2} = 7/6$ are torsion-minimizing structural rationals (Prop.~\ref{prop:G}). The value $g_{K3} = 64/77$ is a \emph{rational approximation} to the arithmetic mean of the four K3 eigenvalues, accurate to $4 \times 10^{-3}$ but not an exact identity: imposing $g_{ss} = 19/6$, $g_{T^2} = 7/6$, $\det(g) = 65/32$ exactly and assuming K3 isotropy forces $g_{K3}^{\mathrm{exact}} = (1755/3724)^{1/4} = 0.828546\ldots$, which does not match any simple closed form in the TCS basis $\{1, \pi, \sqrt{n} : n \leq 110\}$ at present precision. The mismatch reflects intrinsic K3 anisotropy, not a failure of the topological ansatz.

\textbf{K3 near-roundness.} At $s = 0.5$ the four K3 eigenvalues $(0.82209, 0.82770, 0.82974, 0.83167)$ cluster within spread $(\lambda_{\max} - \lambda_{\min})/\bar\lambda = 1.16\%$. The rational approximation $64/77 \approx 0.8312$ matches the arithmetic mean to $4.1 \times 10^{-3}$.

\textbf{Interval-certified measurements.} Interval arithmetic (mpmath.iv, dps $= 50$) in the companion notebook certifies: (i) $\det(g(0.5)) = 65/32$ to better than $10^{-12}$; (ii) each K3 block eigenvalue interval has width $\sim 10^{-12}$, separated by gaps of $O(10^{-3})$ via a Weyl--Bauer--Fike halo; (iii) PSLQ is robustly null for $g_{K3}^{\mathrm{exact}}$ across 15 configurations.

\textbf{1-parameter signature of the K3 block.} An \emph{extended} (post-certificate) Joyce NK run on the K3 block, 9 iterations with cumulative torsion reduction factor $18{,}837$, going beyond the 5-iteration $\sim 3{,}000\times$ reduction of the main certificate (§\ref{sec:torsion}, Table~\ref{tab:torsion_iter}), leaves the normalized K3 eigenvalue-deviation ratios $r_i = (\lambda_i - \bar\lambda)/(\lambda_{\max} - \bar\lambda)$ stable to 15 digits:
\begin{equation}
(r_0, r_1, r_2, r_3) = (-1.47620631853764,\; -0.02477659401258,\; +0.50098291255022,\; +1)
\end{equation}
with scale $\sigma = 3.82755598 \times 10^{-3}$. The naive pattern $(-3/2, 0, 1/2, 1)$ is falsified at $10^{-12}$ precision. The deviation satisfies $\mathrm{dev}_0 + \mathrm{dev}_1 \approx 0$, $\mathrm{dev}_2 \approx 0$, giving the \textbf{1-parameter signature}:
\begin{equation}
r = (-3/2 + \delta,\; -\delta,\; 1/2,\; 1), \qquad \delta \approx 0.02379.
\end{equation}

\textbf{Classification of structural relations:}
\begin{table}[ht]
\centering
\small
\begin{tabular}{p{4.2cm} p{2.8cm} p{8.4cm}}
\toprule
Identity & Status & Justification \\
\midrule
$g_{ss} = 19/6$ & \textbf{Observed, validated} & Torsion minimizer (Prop.~\ref{prop:G}); Lean-checked algebraic consistency \\
$g_{T^2} = 7/6$ & \textbf{Observed, validated} & Torsion minimizer (Prop.~\ref{prop:G}); Lean-checked \\
$g_{K3} \approx 64/77$ & \textbf{Observed, conjectural} & Arithmetic mean to $0.41\%$; exact $g_{K3} = (1755/3724)^{1/4}$ irrational in present basis \\
$\det(g) = 65/32$ & \textbf{Algebraic} & Imposed exactly by construction (§\ref{sec:cheb}); interval-certified to width $1.4 \times 10^{-12}$ \\
$\gamma^2 = 24\pi^2/7 = 4\pi^2/g_{T^2}$ & \textbf{Derived} & $H^1(K3) = 0$ + $T^2$ Hodge Laplacian; NK rate $5.811$ is $0.1\%$ below \\
Naive pattern $(-3/2, 0, 1/2, 1)$ & \textbf{Falsified} & Interval cert rejects at $10^{-12}$; 1-parameter $\delta \approx 0.024$ survives \\
Even Chebyshev modes $= 0$ & \textbf{Certified} & Machine-precision vanishing ($< 10^{-7}$) \\
Eigenvalue constancy & \textbf{Certified} & Drift $\leq 0.18\%$ along neck \\
\bottomrule
\end{tabular}
\caption{Classification of structural identities.}
\end{table}

\subsection{Adiabatic decomposition of the torsion-free metric}

\begin{proposition}[Adiabatic decomposition]\label{prop:E}
Let $g^*$ be the torsion-free $\Gtwo$ metric from Proposition~\ref{prop:A}, with $\|g^* - g_5\|/\|g_5\| \leq 1.35 \times 10^{-7}$. Then $g^* = g_{\mathrm{product}} + E$ where:

\emph{(i)} $\|E\|_F / \|g^*\|_F \leq 2.0 \times 10^{-3}$.

\emph{(ii)} The Hodge Laplacian decomposes as $\Delta(g^*) = \Delta_{\mathrm{seam}} \otimes 1 + 1 \otimes \Delta_{T^2} + 1 \otimes \Delta_{K3} + \delta\Delta$ with $\|\delta\Delta\|/\|\Delta\| \leq 7.8 \times 10^{-3}$.

\emph{(iii)} All eigenvalues satisfy $|\delta\lambda_n|/\lambda_n \leq \kappa(g) \times \varepsilon_{\mathrm{ad}} \leq 0.78\%$, where $\kappa(g) \leq 3.88$ and $\varepsilon_{\mathrm{ad}} = 2.0 \times 10^{-3}$.
\end{proposition}

The gap ratio $\lambda_{22}/\lambda_{21} \approx 14{,}635$ for $\Delta_2$ is stable under this $0.78\%$ perturbation, confirming $b_2 = 21$ robustly.

% ============================================
\section{Torsion Analysis}\label{sec:torsion}
% ============================================

\subsection{Torsion norms and $\Gtwo$ representation theory}

The torsion of a $\Gtwo$-structure $(\varphi, \psi)$ is $T = (d\varphi, d\star\varphi)$. In the seam sector, the nonvanishing contributions come from $\partial_s \varphi$ and $\partial_s \psi$:
\begin{equation}
|d\varphi|^2(s) = 4\, g^{00}(s) \cdot g^{aa'}(s)\, g^{bb'}(s)\, g^{cc'}(s) \cdot (\partial_s \varphi_{abc})(\partial_s \varphi_{a'b'c'}).
\end{equation}

In the Fern\'andez--Gray classification~\cite{FernandezGray1982}, the torsion lies predominantly in the $W_3$ component, with:
\begin{equation}
|d\varphi|^2 / |d\star\varphi|^2 = 1/5 \quad \text{(exact by representation theory)}
\end{equation}
verified to $10^{-10}$ precision at every evaluation point.

\textbf{C\textsuperscript{0} (supremum) norm:}
$\|T\|_{C^0} := \sup_{s \in [-2,3]} \sqrt{|d\varphi|^2(s) + |d\star\varphi|^2(s)}$
evaluated at $N = 100$ Chebyshev collocation nodes on $[0, 1]$. The torsion is a polynomial of degree $\leq 2K = 10$; evaluation at Clenshaw--Curtis nodes is exact for this degree.

\subsection{K3 fiber verification}\label{sec:frech}

The seam sector treats the K3 fiber metric as spatially homogeneous. The extension to the full 7D metric proceeds via an analytical Fr\'echet bound.

\textbf{Step 1: Analytical Fr\'echet bound.}
\begin{equation}
\left\|\frac{\partial T}{\partial g_{K3}}\right\| \leq C_{\mathrm{Fano}} \cdot \|g^{-1}\|^2 \cdot \left\|\frac{dL}{ds}\right\|_F \leq 42 \times 1.224^2 \times 0.014 = 0.881
\end{equation}

\begin{table}[ht]
\centering
\begin{tabular}{ll}
\toprule
Quantity & Value \\
\midrule
Seam-sector torsion $\|T\|_{C^0}^{\mathrm{seam}}$ & $2.984 \times 10^{-5}$ \\
Linear K3 correction to torsion & $\leq 1.59 \times 10^{-3}$ \\
Quadratic K3 correction & $4.95 \times 10^{-12}$ \\
K3-inclusive certified torsion $\|T\|_{C^0}^{7D}$ & $\leq 1.59 \times 10^{-3}$ \\
\bottomrule
\end{tabular}
\caption{Fr\'echet perturbation chain for K3 fiber contribution.}
\end{table}

\textbf{Independent NK certification of the K3 fiber.} The $\CIvvv$ surface admits an independent Newton--Kantorovich certificate via the Donaldson algebraic section method \cite{Donaldson2009} at degree $k=4$ ($126$ sections, $31{,}752$ parameters). Computer-assisted period and monodromy analyses on K3 toric hypersurfaces in a related interval-arithmetic framework appear in \cite{IshigeTakayasu2025}. On a held-out test set of $1{,}000$ fresh points, the Monge--Amp\`ere residual is $\eta_{L^2} = 1.60 \times 10^{-2}$. Two structurally independent $\beta$ sources certify $h < 1/2$: a graph-Laplacian bound $\beta_{\mathrm{Lap}} = 5.66$ gives $h = 7.83 \times 10^{-2}$ (margin $\times 6.4$); a Jacobian pseudoinverse bound $\beta_{\mathrm{Jac}} = 2.25$ at $k=3$ gives $h = 0.188$ (margin $\times 2.7$). The condition fails at $k=2$ ($h = 1.55 > 1/2$), demonstrating that $h < 1/2$ is a genuinely discriminating condition on approximation quality. Full details are in [C].

\subsection{Torsion at the final iterate}

\begin{table}[ht]
\centering
\begin{tabular}{lllll}
\toprule
Norm & $\|d\varphi\|$ & $\|d\star\varphi\|$ & $\|T\|$ & $|d\varphi|^2/|d\star\varphi|^2$ \\
\midrule
$C^0$ & $1.218 \times 10^{-5}$ & $2.724 \times 10^{-5}$ & $2.984 \times 10^{-5}$ & $0.200000$ \\
\bottomrule
\end{tabular}
\caption{Torsion at the final iterate $g_5$. The reported $\|T\| = 2.984 \times 10^{-5}$ is $\sqrt{(\sup|d\varphi|)^2 + (\sup|d\star\varphi|)^2}$ (product of component sups); the direct sup of $\sqrt{|d\varphi|^2 + |d\star\varphi|^2}$ is the slightly tighter $2.949 \times 10^{-5}$ used as $\eta$ in §\ref{sec:nk}.}
\label{tab:torsion_iter}
\end{table}

% ============================================
\section{Newton--Kantorovich Certification}\label{sec:nk}
% ============================================

\subsection{Banach space framework}

\begin{definition}\label{def:banach}
The solution space is the finite-dimensional Banach space $X = \RR^{168}$, parametrizing the Chebyshev--Cholesky coefficients $c = (c_{kj})$. Each $c \in X$ determines a metric $g(s; c) = L(s; c) L(s; c)^T$ via the reconstruction algorithm (§\ref{sec:cheb}). The torsion $T(g(s; c))$ is a polynomial of degree $\leq 2K = 10$ in the Chebyshev variable; the Chebyshev tail beyond degree $K$ is \textbf{identically zero}.
\end{definition}

\begin{definition}\label{def:torsion_map}
The \emph{torsion map} $F: X \to Y = C^0(I, \Lambda^4 \oplus \Lambda^5)$ is defined by $F(g) = (d\varphi[g],\, d\star\varphi[g])$.
\end{definition}

\begin{remark}
Since $L(s; c)$ is a polynomial of degree $K = 5$ in $s$, the torsion $T(g(s; c))$ is a polynomial of degree $\leq 2K = 10$ in the Chebyshev variable. The Chebyshev tail beyond degree $K$ is \emph{identically zero}, there is no truncation error. This distinguishes the present problem from the infinite-dimensional setting of, e.g., the Nirenberg problem on $S^2$, and from a full-7D NK certificate (as pursued in Platt's programme for general special holonomy metrics~\cite{Platt2026}), where the entire metric tensor would be simultaneously discretized.
\end{remark}

\subsection{Abstract NK theorem}

\begin{theorem}[Newton--Kantorovich {\cite{Kantorovich1948,OrtegaRheinboldt1970}}]\label{thm:NK}
Let $X, Y$ be Banach spaces, $U \subset X$ open, and $F: U \to Y$ continuously Fr\'echet differentiable. Suppose $g_0 \in U$ and there exist $\beta, \eta, \omega > 0$ such that:
\begin{enumerate}
\item[(NK1)] $\|F'(g_0)^{-1}\| \leq \beta$.
\item[(NK2)] $\|F(g_0)\| \leq \eta$.
\item[(NK3)] $\|F'(g_1) - F'(g_2)\| \leq \omega \|g_1 - g_2\|$ for all $g_1, g_2 \in U$.
\end{enumerate}
If $h = \beta\eta\omega < 1/2$, then $F$ has a unique zero $g^* \in B(g_0, r_-)$ with
\[
r_- = \frac{1 - \sqrt{1 - 2h}}{\omega}, \qquad r_+ = \frac{1 + \sqrt{1 - 2h}}{\omega},
\]
and convergence is quadratic.
\end{theorem}

\subsection{Invertibility of the linearized operator (NK1)}

\begin{lemma}[Analytical inverse bound]\label{lem:beta}
The linearized operator satisfies $\|F'(c_0)^{-1}\| \leq \beta_a = 0.321$, certified analytically.
\end{lemma}

\begin{proof}
The fundamental mode of the Sturm--Liouville operator $-d/ds(p(s)\, du/ds) = \lambda w(s) u$ on $[0,1]$ has eigenvalue $\lambda_1 \geq \pi^2/g_{00}^{\max}$. The weight $w(s) = \sqrt{\det g} = \sqrt{65/32}$ is constant (algebraically exact). Chebyshev evaluation at Clenshaw--Curtis nodes with Bernstein--Markov inter-node correction gives $g_{00}^{\max} \leq 3.166$. Therefore $\lambda_1 \geq \pi^2/3.166 \geq 3.115$, giving $\beta_a = 1/\lambda_1 \leq 0.321$. \qed
\end{proof}

\subsection{Torsion residual (NK2)}

\begin{lemma}[Torsion residual]\label{lem:eta}
$\|F(g_0)\| = \sup_{s \in [0,1]} \sqrt{|d\varphi|^2(s) + |d\star\varphi|^2(s)} \leq 2.949 \times 10^{-5}$, evaluated directly at Clenshaw--Curtis collocation nodes (exact for the degree-10 torsion polynomial). The slightly larger figure $2.984 \times 10^{-5}$ in Table~\ref{tab:torsion_iter} is the conservative product $\sqrt{(\sup |d\varphi|)^2 + (\sup |d\star\varphi|)^2}$ of the two component sups; the direct-sup value $2.949 \times 10^{-5}$ is the one used as $\eta$ in the NK certificate.
\end{lemma}

\subsection{Lipschitz constant (NK3)}

\begin{lemma}[Lipschitz bound]\label{lem:omega}
$\|F'(g_1) - F'(g_2)\| \leq \omega = 1.636 \times 10^{-3}$ for all $g_1, g_2 \in B(g_0, r_+)$. Certified via Chebyshev derivative bounds (Bernstein--Markov inequality on $D^2T$) with no finite differences.
\end{lemma}

\subsection{Assembly of the certificate}

\begin{proposition}[NK certificate]\label{prop:nkcert}
Both rows satisfy $h = \beta\eta\omega < 1/2$:
\begin{center}
\begin{tabular}{lllll}
\toprule
Certificate & $\beta$ & $\eta$ & $\omega$ & $h = \beta\eta\omega$ \\
\midrule
Analytical & $0.321$ & $2.949 \times 10^{-5}$ & $9.47 \times 10^{-4}$ & $\mathbf{8.95 \times 10^{-9}}$ \\
Numerical (sharper $h$) & $0.02961$ & $2.949 \times 10^{-5}$ & $1.636 \times 10^{-3}$ & $\mathbf{1.43 \times 10^{-9}}$ \\
\bottomrule
\end{tabular}
\end{center}
\end{proposition}

The maximum $\omega$ still permitting certification is $\omega_{\max} = 1/(2\beta\eta) = 5.28 \times 10^4$, a factor $5.6 \times 10^7$ above the computed value.

\subsection{Existence, uniqueness, and holonomy}\label{sec:holonomy}

\textbf{Proposition~\ref{prop:A}} follows from Proposition~\ref{prop:nkcert} and Theorem~\ref{thm:NK}:
\begin{equation}
r_- = 8.73 \times 10^{-7}, \qquad r_+ = 1{,}222, \qquad \delta g/g \leq 1.35 \times 10^{-7}.
\end{equation}

\begin{figure}[ht]
\centering
\includegraphics[width=\textwidth]{fig_chebyshev_nk_paper_a.png}
\caption{Chebyshev mode amplitudes and NK certificate. \emph{(a)} Per-mode RMS magnitudes $|c_k|_{\mathrm{rms}}$ (log scale) over the 28 Cholesky entries: $k = 0$ carries $> 99.99\%$ of the signal; even modes $k \in \{2, 4\}$ vanish at machine precision (central antisymmetry of the Cholesky factor); odd modes $k \in \{1, 3, 5\}$ provide small TCS-matching corrections at the $10^{-3}$/$10^{-5}$ level. \emph{(b)} NK certificate schematic showing the contraction quotient $h = \beta\eta\omega$ assembled from the analytical inverse bound, the certified torsion residual at Clenshaw--Curtis nodes, and the Bernstein--Markov Lipschitz bound; the analytical certificate gives $h = 8.95 \times 10^{-9}$ (margin $\times 5.6 \times 10^7$ vs the threshold $1/2$).}
\label{fig:chebyshev_nk}
\end{figure}

\begin{corollary}[Holonomy]
If $M$ is compact and simply connected, then $\mathrm{Hol}(g^*) = \Gtwo$ (simple connectivity gives $b_1(M) = 0$, the hypothesis used by Joyce \cite[Prop.~11.2.3]{Joyce2000}).
\end{corollary}

\textbf{Numerical indicators} (at $g_5$):
$|\varphi|^2 = 42.0000000000$ (error $2.1 \times 10^{-14}$); $\|Rm\|_{C^0} = 4.20 \times 10^{-5}$; $\|\Gamma\|_{C^0} = 2.81 \times 10^{-6}$.

\subsection{From 1D certificate to 7D existence}

\begin{theorem}[1D$\to$7D Reduction]
Let $g_5 = g_{\mathrm{seam}}(s) \oplus g_{T^2} \oplus g_{K3}(y)$ satisfy the hypotheses of Proposition~\ref{prop:A}. Let $C_{\mathrm{red}} \leq 0.881$ (analytical Fr\'echet bound, §\ref{sec:frech}) and $\delta_{\mathrm{fib}} \leq 1.80 \times 10^{-3}$ (K3 fiber variation). Then:
\begin{equation}
\|T(g^*_{7D})\|_{C^0} \leq C_{\mathrm{red}} \cdot \delta_{\mathrm{fib}} \leq 1.59 \times 10^{-3} \ll \varepsilon_0 \approx 0.1 \quad \text{(Joyce threshold)}.
\end{equation}
By Joyce \cite[Thm.~11.6.1]{Joyce2000}, a torsion-free $\Gtwo$ structure exists on any compact extension of $g^*_{7D}$.
\end{theorem}

\begin{table}[ht]
\centering
\small
\begin{tabular}{llll}
\toprule
Hypothesis & Statement & Status \\
\midrule
J1 & $M$ is a compact 7-manifold & \textbf{Conditional} \\
J2 & $M$ carries a $\Gtwo$ structure $\varphi_0$ & \textbf{Certified} \\
J3 & $\|T(\varphi_0)\|_{C^0} < \varepsilon_0$ & \textbf{Certified} ($1.59 \times 10^{-3}$, factor $\times 63$) \\
J4 & $M$ simply connected ($\pi_1 = 1$) & \textbf{Conditional} \\
J5 & $b_1(M) = 0$ & \textbf{Conditional} (follows from J4) \\
J6 & Bounded geometry on neck & \textbf{Certified} \\
\bottomrule
\end{tabular}
\caption{Joyce perturbation theorem hypothesis verification.}
\end{table}

% ============================================
\section{Spectral Democracy}\label{sec:democracy}
% ============================================

\subsection{The Weitzenb\"ock identity on product metrics}

The general Weitzenb\"ock identity for 1-forms: $\Delta_1 \alpha = \nabla^* \nabla \alpha + \mathrm{Ric}(\alpha^\sharp)^\flat$. When $\mathrm{Ric}(g) = 0$, this simplifies to $\Delta_1 = \nabla^* \nabla$.

\subsection{Proof of Theorem~\ref{thm:D}}

We prove that for $\alpha = f(s)\,d\theta$ with $g_{\theta\theta}$ constant, $\Delta_1(f\,d\theta) = (\Delta_0 f)\,d\theta$.

\textbf{Step 1: Christoffel symbols.} For the product metric, the only nonvanishing Christoffel symbols involving the $s$-direction are $\Gamma^s_{ss} = g_{ss}'/(2g_{ss})$ and $\Gamma^a_{sa} = g_{aa}'/(2g_{aa})$ for $a \neq s$. Since $g_{\theta\theta}$ is constant, $\Gamma^s_{\theta\theta} = 0$ and $\Gamma^\theta_{s\theta} = 0$.

\textbf{Step 2: Covariant derivatives of $\alpha = f(s)\,d\theta$.} We compute $\nabla_s \alpha = (\partial_s f)\,d\theta$ (since $\Gamma^\theta_{s\theta} = 0$). For fiber directions: $\nabla_a \alpha = 0$. For the $\theta$-direction: $\nabla_\theta \alpha = 0$.

\textbf{Step 3: Rough Laplacian.} The only nonzero covariant derivative is $\nabla_s \alpha = f'(s)\,d\theta$. The rough Laplacian gives $\nabla^*\nabla \alpha = -g^{ss}(f'' - \Gamma^s_{ss} f')\,d\theta$. The scalar Laplacian on functions of $s$ alone is:
\[
\Delta_0 f = -g^{ss}\bigl(f'' + (\log\sqrt{\det g})'\, f'\bigr).
\]
The difference between these two expressions is the single term $g'_{\theta\theta}/(2g_{\theta\theta})$. Since $g_{\theta\theta}$ is constant by hypothesis, this term vanishes, giving $\Delta_1(f\,d\theta) = (\Delta_0 f)\,d\theta$. The Weitzenb\"ock curvature term vanishes by Ricci-flatness. \qed

\begin{remark}
The constancy of $g_{\theta\theta}$ is essential. For our metric, the $T^2$ fiber is isotropic to $3 \times 10^{-7}$, so corrections are bounded at the $10^{-7}$ level. Numerical verification on the certified metric is presented in the companion paper [B].
\end{remark}

% ============================================
\section{Discussion}\label{sec:discussion}
% ============================================

\subsection{Summary}

We have presented:
\begin{enumerate}
\item A computer-assisted certification of a torsion-free $\Gtwo$ metric on a neck model, using interval arithmetic with zero finite differences (§§\ref{sec:metric}--\ref{sec:nk}).
\item A complete metric decomposition reducing 169 parameters to 5 structural constants (§\ref{sec:decomp}).
\item An analytical spectral democracy theorem for product-type $\Gtwo$ metrics (§\ref{sec:democracy}).
\item A TCS reduction theorem extending the 1D certificate to 7D via analytical Fr\'echet bound and Joyce perturbation (§\ref{sec:nk}).
\end{enumerate}

\subsection{Comparison with existing work}

Concurrent work by Heyes--Hirst--S\'a Earp--Silva \cite{HeyesHirst2026} applies neural networks to approximate $\Gtwo$-structures on contact Calabi--Yau 7-manifolds. Their approach and ours are complementary: \cite{HeyesHirst2026} works on non-compact geometries with neural architectures, while the present paper works on a TCS-type neck model with analytical parametrization and a convergence certificate. Extension of either result to a closed compact $\Gtwo$-holonomy manifold remains conditional on a compatible compact construction.

\subsection{Limitations}

\textbf{Topological realization.} The pair $(b_2, b_3) = (21, 77)$ does not appear among the 252 Joyce orbifold types or CHNP TCS examples. Orthogonal TCS is excluded by parity. Non-orthogonal TCS and other constructions remain open.

\textbf{Certification method.} The NK certificate uses interval arithmetic (mpmath.iv, 50-digit precision) with zero finite differences. The structural constants are formally verified in Lean~4 (4 axioms, zero incomplete proofs). The certificate has not been fully formalized in a proof assistant; recent progress on Mathlib elliptic-PDE infrastructure \cite{DGNMLean2026} suggests this becomes increasingly tractable.

\subsection{Open questions}

\begin{enumerate}
\item \textbf{Geometric realization of $(21, 77)$.} Does there exist a compact $\Gtwo$ manifold with these Betti numbers? Orthogonal TCS is excluded; other constructions remain open.
\item \textbf{Proof-assistant formalization.} Can the NK certificate be fully formalized in Lean~4?
\item \textbf{Fiber-coupled corrections.} What is the spectral shift from explicit K3 mode coupling beyond the adiabatic approximation?
\end{enumerate}

% ============================================
\section*{Appendix A. Metric Data}
\addcontentsline{toc}{section}{Appendix A. Metric Data}
% ============================================

The 169 parameters will be deposited at Zenodo upon publication. Reconstruction follows the algorithm in §\ref{sec:cheb}. The K3 metric $g_{K3}(y)$ is computed via the cymyc neural network \cite{LarforsLukasRuehle2022} on $\CIvvvv \subset \Psix$ with $220{,}000$ sample points.

% ============================================
\section*{Appendix B. NK Proof Chain}
\addcontentsline{toc}{section}{Appendix B. NK Proof Chain}
% ============================================

\textbf{Unified certificate} (12 steps):
\begin{enumerate}
\item Chebyshev coefficient bounds on metric entries (triangle inequality).
\item Softplus structural positivity: $L_{ii} > 0$ for all $s$.
\item SPD from Cholesky: $g = LL^T$ with $L_{ii} > 0$.
\item Algebraic determinant: $\det(g) = 65/32$.
\item Gershgorin eigenvalue enclosure: $\lambda_{\min} \geq 0.817$.
\item Torsion at Clenshaw--Curtis collocation nodes (exact for degree $\leq 10$).
\item K3 Fr\'echet derivative over $220{,}000$ fiber points (§\ref{sec:frech}).
\item K3 fiber correction: $\delta_{K3} \leq 1.59 \times 10^{-3}$ (analytical Fr\'echet, zero FD).
\item Lipschitz bound: $\omega \leq 9.47 \times 10^{-4}$ (Bernstein--Markov, zero FD).
\item Operator inverse bound: $\beta = 1/33.77 = 0.02961$ (analytical: $\beta_a = 0.321$).
\item NK contraction: $h = 8.95 \times 10^{-9} < 0.5$ (analytical $\beta$).
\item TCS reduction: 7D torsion $\leq \delta_{K3} \leq 1.59 \times 10^{-3} < \varepsilon_0 = 0.1$ (Joyce). \qed
\end{enumerate}

All bounds verified at 50-digit precision via mpmath. Companion interval-arithmetic notebooks certify the metric-decomposition observations independently.

% ============================================
\section*{Appendix C. K3 Fiber Clarification}
\addcontentsline{toc}{section}{Appendix C. K3 Fiber Clarification}
% ============================================

\begin{table}[ht]
\centering
\begin{tabular}{ll}
\toprule
Parameter & Value \\
\midrule
Sample points & $220{,}000$ \\
Metric accuracy $\sigma$ & $0.011$ \\
SPD fraction & $100\%$ ($220{,}000/220{,}000$) \\
Hyperk\"ahler triple & $J_I \cdot J_J = J_K$ to $2.2 \times 10^{-10}$ \\
Intrinsic K3 variation & $69.8\%$ ($\sigma/\|g\|$) \\
NK-scaled variation in $\Gtwo$ & $0.037\%$ \\
\bottomrule
\end{tabular}
\caption{K3 fiber metric parameters. The ``NK-scaled variation'' (0.037\%) is the relevant measure for the product-type construction; the ``intrinsic variation'' (69.8\%) reflects the K3 Riemann curvature and is expected.}
\end{table}

% ============================================
\section*{Author's note}
% ============================================

This framework was developed through sustained collaboration between the author and several AI systems, primarily Claude (Anthropic), with contributions from GPT (OpenAI), Gemini (Google), and Aristotle (Harmonic) for specific mathematical insights. Architectural decisions and many key derivations emerged from iterative dialogue sessions. This collaboration follows a transparent crediting approach for AI-assisted mathematical research.

% ============================================
\section*{Disclosure statement}
% ============================================

The author reports there are no competing interests to declare.

% ============================================
\section*{Author's related works}
% ============================================

\begin{itemize}
\item[{[B]}] de La Fourni\`ere, B. \emph{Spectral Geometry of an Explicit $\Gtwo$ Metric: Laplacian Spectrum and Harmonic Forms}. Zenodo, 2026. DOI: \href{https://doi.org/10.5281/zenodo.19893371}{10.5281/zenodo.19893371}.
\item[{[C]}] de La Fourni\`ere, B. \emph{Newton--Kantorovich diagnostics on a Donaldson K3 metric: two $\beta$ estimates, machine-checked inequalities}. Zenodo, 2026. DOI: \href{https://doi.org/10.5281/zenodo.19708916}{10.5281/zenodo.19708916}.
\end{itemize}

% ============================================
\bibliographystyle{plain}
\bibliography{g2_certified_neck}
% ============================================

\end{document}
